\section{Bounded carrier theorem for two disabled mixed
tops}

\textbf{Proof-first theorem, 2026-08-12 PDT.} This note closes the
one-active single-linkage face with two disabled top complexes. It is
independent of the separated-scale operator and uses no support or
orientation enumeration.

\subsection{1. Setting and theorem}

Put

\[
 q_B=A+B,\qquad q_C=A+C,\qquad
 {\cal P}=\{B,C,2B,2C,B+C\}.                                 \tag{1.1}
\]

Fix one strongly connected linkage support

\[
 \{q_B,q_C\}\subseteq{\cal C}
      \subseteq\{0,q_B,q_C\}\cup{\cal P},                    \tag{1.2}
\]

with arbitrary positive labelled rates. The no-fast entrance is

\[
                              x_n=(n,0,0).                     \tag{1.3}
\]

For fixed \(\ell\), put

\[
 G_\ell(x)=K_\ell+\sum_{i=A,B,C}\log(x_i!)+\ell\cdot x\ge1,
 \qquad W_\ell=G_\ell^4.                                     \tag{1.4}
\]

\begin{quote}
\textbf{Theorem 1.1 (bounded two-top service).} In one fixed closed
irreducible class, exactly one of the following applies.

\begin{enumerate}
\def\labelenumi{\arabic{enumi}.}
\tightlist
\item
  If \(0\notin{\cal C}\), (1.3) is absorbing and its class is a
  singleton.
\item
  If \(0\in{\cal C}\) and \({\cal C}\cap{\cal P}=\varnothing\), then
  \(A-B-C\) is invariant. The cofactor-free part of a fixed class
  contains at most one state (1.3).
\item
  If \(0\in{\cal C}\) and \({\cal C}\cap{\cal P}\ne\varnothing\), there
  is an included all-reaction stopping time \(\tau_n\) such that, for
  every fixed \(p\), \[
  \begin{aligned}
   \mathbb P(E_n)&\le C/n,\\
   \mathbb E[(1+|X_{\tau_n}-x_n|+\tau_n)^p;E_n]&\le C_p/n,\\
   \mathbb E(1+\tau_n)^p&\le C_p,\\
   A_{\tau_n}&=n-1,\quad B_{\tau_n}+C_{\tau_n}\le3
           \quad\hbox{on }E_n^c,
  \end{aligned}
  \tag{1.5}
  \] and, for all large \(n\), \[
   \mathbb E_{x_n}[
     W_\ell(X_{\tau_n})-W_\ell(x_n)+\tau_n]
      \le-cG_\ell(x_n)^3\log n.                            \tag{1.6}
  \]
\end{enumerate}

Here \(E_n\) is the included first lower-source competitor during a
mixed-top carrier phase. All physical clocks are retained.
\end{quote}

Constants may depend on the graph, rates, class, and \(\ell\), but not
on \(n\).

\subsection{2. Frozen and invariant
alternatives}

The complex weights of

\[
                              H=A-B-C                         \tag{2.1}
\]

are zero on \(0,q_B,q_C\), \(-1\) on \(B,C\), and \(-2\) on
\(2B,2C,B+C\).

If \(0\) is absent, no source is enabled at (1.3), so the state is
absorbing. If \({\cal P}\) is absent, every reaction has zero
\(H\)-reward. At \(B=C=0\), the invariant value is \(H=A\), so a fixed
class contains at most one cofactor-free value of \(A\). This proves
alternatives 1--2.

\subsection{3. The clean geometric attempt
kernel}

Assume alternative 3. At the no-fast face, \(0\) is the only enabled
source, with aggregate rate

\[
                 \lambda_0=\sum_{0\to z}\kappa_{0z}>0.        \tag{3.1}
\]

Choose a simple directed path from \(0\) to its first vertex in
\({\cal P}\). Before its terminal vertex, it lies in \(\{0,q_B,q_C\}\).
Simplicity prevents an intermediate return to \(0\). After the launch,
every prescribed source is \(q_B\) or \(q_C\), and is physically enabled
by the preceding target with the same residual population.

Temporarily suppress lower-source clocks during a top phase. A carrier
moves on the finite transient graph \(\{q_B,q_C\}\), absorbed at
\(\{0\}\cup{\cal P}\). Conditional outgoing probabilities at a top
source are independent of \(n\), because the outgoing clocks have their
common falling-factorial source. The chosen path therefore gives a fixed

\[
                              p_*>0                            \tag{3.2}
\]

chance that an attempt reaches \({\cal P}\). Its internal top-event
count is phase type and its physical duration has moments \(O(n^{-p})\).

Absorption at \(0\) returns to the exact population \((n,0,0)\). Repeat
attempts. Their count is dominated by a geometric variable of parameter
\(p_*\). Each no-fast visit waits an exponential time of rate
\(\lambda_0\). Thus the clean attempt count and total duration have
every fixed moment.

\subsection{4. Restoring every physical
clock}

During any pre-success top phase, the inactive population is at most
one. Every lower source, including \(0\) and any enabled member of
\({\cal P}\), therefore has aggregate rate bounded by a graph-dependent
constant. The aggregate enabled top rate is at least \(cn\). Hence

\[
 \mathbb P\{\hbox{a lower firing before clean top absorption}\}
       \le C/n                                                \tag{4.1}
\]

per attempt. Stop at and include the first such firing.

At a clean \({\cal P}\)-target, the inactive population has molecularity
one or two. At least one of \(q_B,q_C\) is enabled, and its rate is at
least \(cn\). Follow top reactions until the first top-to-lower exit.
Top-to-top conversions preserve \(A\); this exit lowers \(A\) by one.
Before it, inactive population is at most two and the lower-source
aggregate rate is bounded. Its included competitor probability is again
\(O(n^{-1})\).

Sum these races over the geometric attempt count. The causing lower jump
has a binary target, so its actual endpoint is bounded relative to the
clean carrier state. Geometric-sum moment formulas give the first three
lines of (1.5).

On a clean successful attempt, a launch \(0\to q_B\) or \(0\to q_C\)
raises \(A\) by one and its first top exit cancels that entry. Reaching
\({\cal P}\) leaves an inactive carrier, whose next top exit is
unpaired. A direct \(0\to{\cal P}\) launch supplies that unpaired exit
immediately. Thus

\[
                  A_{\tau_n}=n-1,\qquad
                  B_{\tau_n}+C_{\tau_n}\le3.                  \tag{4.2}
\]

This proves all of (1.5). No lower clock has been erased: it is either a
designated no-fast launch or the included event \(E_n\).

\subsection{5. Entropy and common fourth
power}

On the clean event, (4.2) gives

\[
                    G_\ell(X_{\tau_n})-G_\ell(x_n)
                        \le-\log n+C_\ell.                    \tag{5.1}
\]

On \(E_n\), the centered displacement before the causing jump has
geometric moments, and the causing jump is bounded. Consequently

\[
 \mathbb E[
   |G_\ell(X_{\tau_n})-G_\ell(x_n)|^p;E_n]
       \le {C_{p,\ell}(\log(n+e))^p\over n}.                  \tag{5.2}
\]

Together with the duration moments, this yields

\[
 \mathbb E[
   G_\ell(X_{\tau_n})-G_\ell(x_n)+\tau_n]
       \le-\tfrac12\log n                                    \tag{5.3}
\]

for all large \(n\).

Since \(G_\ell(x_n)\asymp n\log n\), the exact identity

\[
 (G+z)^4-G^4=4G^3z+6G^2z^2+4Gz^3+z^4                       \tag{5.4}
\]

and (5.2) show that the last three expected terms are
\(O(G_\ell(x_n)^2(\log n)^2)\), lower order than
\(G_\ell(x_n)^3\log n\). The physical-duration term is lower order as
well. This proves (1.6). \(\square\)

\subsection{6. Scope}

This theorem treats exactly the bounded-base one-active shape with both
\(A+B\) and \(A+C\) disabled at \((n,0,0)\). It is independent of the
mesoscopic separated theorem. It proves arbitrary orientations and rates
from a simple strong-graph path, a finite phase-type carrier, and an
included all-clock race; it uses no finite enumeration.
